\diarytitle{0125}{保存場の判定とポテンシャル関数による線積分の計算}

ベクトル場
\[
\bm{F}(x,y) = \left(2xy + \cos x,\; x^{2} + 2y\right)
\]
について、以下の問に答えよ。
\begin{enumerate}
\item $\bm{F}$ が保存場（$\bm{F} = \nabla \phi$ となるポテンシャル $\phi$ が存在する場）であることを、$xy$ 平面全体で $\partial P/\partial y = \partial Q/\partial x$ が成り立つことを示すことで確認せよ。
\item ポテンシャル関数 $\phi(x,y)$（$\bm{F} = \nabla \phi$、すなわち $\phi_x = P$, $\phi_y = Q$）を求めよ。
\item 経路によらないことを用いて、点 $(0,0)$ から点 $(1,\pi)$ までの任意の滑らかな曲線 $C$ に沿う線積分 $\displaystyle\int_{C} \bm{F}\cdot d\bm{r}$ を求めよ。
\end{enumerate}
